\section{Testing Results}

\subsection{Environment and Methodology}

Describe the hardware, software, datasets, and testing procedure.

\subsection{Test Cases}

For each test, record the concrete input or scenario under \emph{Test Case},
and classify its role under \emph{Test Type}, such as a functional test, a
medium-scale data test, or a large-scale data test. Also record the expected
result, actual result, and status. For a failed test, explain the possible
cause and correction in the discussion below the table.

\begin{table}[htbp]
  \centering
  \caption{Testing results}
  \label{tab:test-results}
  \small
  \begin{tabularx}{\textwidth}{c p{0.20\textwidth} p{0.18\textwidth} X X c}
    \toprule
    ID & Test Case & Test Type & Expected Result & Actual Result & Status \\
    \midrule
    1 & Small hand-checkable input & Functional test & Correct output & Same as expected & Pass \\
    2 & $10^4$ randomly generated records & Medium-scale data test & Correct output within the target time & Record the measured result & Pass \\
    3 & $10^6$ randomly generated records & Large-scale data test & Correct output without failure & Record the measured result & Pass \\
    \bottomrule
  \end{tabularx}
\end{table}

\subsection{Results and Discussion}

Summarize what the tests demonstrate and discuss unexpected behavior or
limitations.
